\chapter{Progettazione e codifica}
\label{chap:progettazione-codifica}

Il progetto \textit{ErpAI} è basato su un'architettura a componenti, nella quale il backend principale è realizzato tramite Laravel ed è responsabile della gestione delle richieste, 
della logica applicativa e dell'accesso ai dati. L'applicazione comunica con i servizi di persistenza e supporto, tra cui MariaDB e Redis, mentre ulteriori componenti esterni 
sono utilizzati per attività specifiche, quali il forecasting e l'ottimizzazione della produzione.

L'attività di stage non ha previsto una modifica dell'architettura generale dell'applicazione, ma l'estensione delle componenti già presenti attraverso l'introduzione del modulo 
XAI e del sistema di alert. Le nuove funzionalità sono state integrate nell'architettura esistente, riutilizzandone i meccanismi di autenticazione, autorizzazione, accesso ai dati 
e gestione delle operazioni asincrone.

Per quanto riguarda l'implementazione dei moduli, è stata adottata una struttura comune basata sulla separazione tra interfaccia, implementazione condivisa e logica specifica delle singole componenti. Tale organizzazione consente di mantenere una struttura omogenea tra i diversi moduli e di favorirne l'estendibilità.

\begin{itemize}
  \item \textbf{Contract:} interfaccia che dichiara i principali metodi che devono essere implementati dalle classi concrete. In questo modo viene definito un contratto 
    comune per le diverse implementazioni, favorendo la compatibilità con le componenti esistenti e l'introduzione di nuove implementazioni.
  \item \textbf{Abstract Class:} classe astratta che implementa il \texttt{Contract} e fornisce l'implementazione dei metodi comuni alle diverse classi concrete, come i metodi 
    getter e setter. Le sottoclassi possono quindi concentrarsi sull'implementazione della logica specifica senza dover replicare il comportamento condiviso.
  \item \textbf{Classi Concrete:} classi che estendono la \texttt{Abstract Class} e implementano la logica specifica della singola componente, differenziandosi in base alle 
    funzionalità che devono fornire.
\end{itemize}


\newpage
\subsection{Struttura del progetto}
\begin{listing}[H]
\begin{minted}{markdown}
erpai/
  ├── app/
  │    ├── Console/Commands/
  |    ├── Contracts/
  |    |     ├── AiChat/    - - - - - - - > AiToolContract.php       
  |    |     ├── Alerts/    - - - - - - - > AlertRuleContract.php 
  |    ├── Filament/
  |    |     ├── Clusters/
  |    |     ├── Pages/
  |    |     └── Resources/
  |    ├── Jobs/            - - - - - - - > AlertJob.php / BroadcastAlertJob.php
  |    ├── Models/ 
  |    ├─- Observers/       - - - - - - - > AuditObserver.php
  |    └── Services/
  |         ├── AiChat/
  |         |      ├── Eval/
  |         |      ├── Prompts/
  |         |      └── Tools/
  |         |      
  |         └── Alerts/
  |
  ├── bootstrap/app.php
  ├── database/
  |     ├── factories/
  |     └── migratons/
  └── tests/alerts/
\end{minted}
\caption{Struttura del progetto}
\label{listing:folders}
\end{listing}

\section{Modulo XAI}
La logica del modulo XAI è organizzata all'interno della cartella \newline\texttt{app/Services/AiChat}. Al suo interno sono raccolte le componenti responsabili della gestione 
dei provider dei modelli linguistici, dei tool, dei prompt e delle funzioni di valutazione (\textit{eval}), queste ultime descritte in dettaglio nella sezione dedicata al 
testing.

\subsection{Tool}
I tool rappresentano il punto di accesso del modello linguistico ai dati e alle funzionalità del gestionale. Attraverso di essi il modello può richiedere informazioni 
al sistema o eseguire operazioni, nel rispetto dei permessi configurati.
\newline
Ogni tool implementa l'interfaccia \texttt{AiToolContract}, che definisce le operazioni comuni che devono essere fornite dalle diverse implementazioni. In particolare, 
il contratto espone i metodi necessari al recupero dei metadati del tool, utilizzati anche per la generazione delle relative etichette nell'interfaccia utente, alla definizione 
dei permessi richiesti e all'esecuzione del tool.
\newline
L'interfaccia viene implementata dalla classe astratta \texttt{AbstractTool}, che fornisce il comportamento condiviso relativo alla gestione e alla verifica dei permessi. Le 
singole classi concrete estendono quindi \texttt{AbstractTool} e implementano la logica specifica necessaria per accedere ai dati o alle funzionalità messe a disposizione dal 
relativo tool.
\newline
Questa struttura permette di mantenere separata la logica comune da quella specifica dei singoli tool, facilitando l'introduzione di nuove funzionalità e mantenendo uniforme 
la gestione delle autorizzazioni.

\subsubsection*{FinderTool}
Per illustrare l'implementazione dei tool viene considerato \texttt{FinderTool}, sviluppato per consentire al modello linguistico di accedere ai dati del gestionale e 
utilizzare informazioni reali durante la generazione delle risposte, riducendo il rischio di produrre informazioni non presenti nel sistema.
\newline
Il tool è composto dai seguenti elementi principali:
\begin{itemize}
  \item \texttt{shortDescription}: descrizione sintetica esposta al modello linguistico, utilizzata per consentirgli di individuare il tool più appropriato in base alla 
    richiesta dell'utente.
  \item \texttt{description}: descrizione dettagliata contenente le istruzioni necessarie per guidare il modello nell'utilizzo corretto del tool.
  \item \texttt{parameters()}: definizione, attraverso uno schema JSON, dei parametri che il modello deve fornire per poter eseguire il tool. Nel caso di \texttt{FinderTool}, 
    lo schema specifica le operazioni da effettuare sui dati e i modelli Eloquent sui quali eseguire le interrogazioni.
  \item \texttt{execute}: metodo responsabile dell'elaborazione dei parametri forniti dal modello e dell'esecuzione dell'operazione richiesta. Nel caso di \texttt{FinderTool}, i 
    parametri definiti attraverso \texttt{parameters()} vengono elaborati per costruire ed eseguire una query sul database, consentendo di recuperare informazioni relative allo 
    stato di una specifica entità del gestionale, come un ordine o una macchina. I risultati ottenuti vengono quindi restituiti al modello linguistico, che li utilizza per 
    elaborare la risposta finale in linguaggio naturale.
\end{itemize}

\subsection{Provider}
I provider rappresentano il livello di integrazione tra il modulo XAI e i modelli linguistici. Essi gestiscono le richieste provenienti dall'applicazione e si occupano di 
comunicare con i servizi che espongono i modelli linguistici attraverso le relative API. In questo modo, la logica specifica necessaria per interagire con ciascun provider 
viene separata dal resto del modulo XAI.
\newline
Ogni provider implementa l'interfaccia \texttt{AiProviderContract}, che definisce le operazioni comuni richieste dalle diverse implementazioni. In particolare, il contratto 
espone i metodi necessari al recupero dei metadati dei modelli, utilizzati anche per la generazione delle relative etichette nell'interfaccia utente, nonché alla configurazione 
delle modalità di \textit{thinking} e di \textit{streaming} e alla gestione della risposta elaborata dal modello linguistico.
\newline
L'interfaccia viene implementata dalla classe astratta \texttt{AbstractAiProvider}, che fornisce il comportamento condiviso tra le diverse implementazioni, occupandosi della 
gestione dei parametri comuni e dell'interazione con i tool. Le singole classi concrete estendono quindi \texttt{AbstractAiProvider} e implementano la logica specifica richiesta 
dai diversi provider e modelli.
\newline
Tra le funzionalità gestite dalle implementazioni concrete rientrano la gestione delle chiamate ai tool, il rispetto dei limiti imposti dal modello relativamente ai tempi di 
elaborazione e al numero di token in ingresso e in uscita, la gestione della modalità di \textit{thinking}, quando supportata dal modello, e la gestione dello \textit{streaming}. 
Quest'ultima permette di trasmettere la risposta progressivamente, suddividendola in più porzioni, riducendo così il tempo percepito dall'utente prima della visualizzazione del 
contenuto.
\newline
Vi è infine il metodo \texttt{reply}, responsabile della gestione dell'interazione con il modello linguistico. Esso si occupa del caricamento dei tool disponibili, della 
costruzione della conversazione e del relativo contesto, nonché della configurazione dei parametri relativi ai limiti e alle modalità di risposta. Una volta predisposti i 
dati necessari, il metodo effettua la chiamata alle API del provider e restituisce la risposta elaborata dal modello.

\subsubsection*{RunpodOllamaProvider}
Per illustrare l'implementazione dei provider viene considerato \newline\texttt{RunpodOllamaProvider}, sviluppato per consentire al modulo XAI di utilizzare un modello linguistico 
eseguito tramite Ollama\footnote{\cite{site:ollama}} su un'istanza Runpod\footnote{\cite{site:runpod}}. Durante lo sviluppo, questo provider è stato quello principalmente utilizzato, grazie alla disponibilità del modello open source 
Qwen3.6:35B\footnote{\cite{site:qwen}} e alle prestazioni di risposta ottenute nell'ambiente di esecuzione utilizzato.
\newline
L'implementazione di \texttt{RunpodOllamaProvider} è volutamente ridotta, poiché la maggior parte della logica comune ai diversi provider è implementata all'interno di 
\texttt{AbstractAiProvider}. La classe concreta si occupa principalmente di fornire i parametri specifici necessari per l'accesso al servizio, quali l'URL dell'API, la chiave 
di autenticazione e i limiti relativi al numero di token supportati dal modello.

\subsection{Services}
I \texttt{Services} costituiscono il livello di coordinamento tra l'interfaccia utente e le componenti del modulo XAI. In questa componente è stata implementata la classe 
\texttt{AiChatService}, responsabile della gestione delle principali operazioni necessarie per l'interazione con il chatbot.

\subsubsection*{AiChatService}
La classe \texttt{AiChatService} coordina l'intero ciclo di utilizzo del modulo XAI. In particolare, gestisce il contesto associato alla richiesta, composto dalle informazioni 
provenienti dalla pagina del gestionale e dalla conversazione corrente. Il contesto permette al modello linguistico di interpretare correttamente anche richieste che, considerate 
singolarmente, potrebbero risultare incomplete.
\newline
La classe si occupa inoltre della verifica dei permessi necessari per l'accesso alla \textit{sidebar} del chatbot e della gestione delle conversazioni, inclusa la creazione di 
nuove conversazioni e il recupero di quelle precedentemente avviate. Un'ulteriore responsabilità riguarda la registrazione e la gestione dei tool disponibili per il provider 
selezionato.
\newline
Infine, \texttt{AiChatService} espone il metodo \texttt{sendMessage}, che coordina l'invio della richiesta al modello linguistico. Il metodo raccoglie il contesto della 
conversazione, il modello configurato, i tool disponibili e il messaggio dell'utente, delegando quindi al provider la gestione della comunicazione con il modello e restituendo 
il risultato alla componente chiamante.


\section{Sistema di generazione degli alert}
La logica del sistema di generazione degli alert è organizzata in diverse componenti, ciascuna con una responsabilità specifica. Il sistema è basato su un insieme di 
regole configurabili che verificano periodicamente lo stato del gestionale e individuano eventuali condizioni che richiedono l'intervento dell'utente.

\subsection{Rules}
Le \textit{rules} rappresentano le componenti responsabili della verifica delle condizioni che possono generare un alert. Esse effettuano interrogazioni al database per 
individuare eventuali violazioni delle regole configurate dall'utente. La maggior parte delle regole prevede inoltre parametri configurabili, che permettono di personalizzare 
le condizioni di attivazione degli alert in base alle esigenze dell'utente.
\newline
Analogamente ai tool del modulo XAI, le regole implementano l'interfaccia \texttt{AlertRuleContract}, che definisce le operazioni comuni che devono essere fornite dalle diverse 
implementazioni. Il contratto viene implementato dalla classe astratta \texttt{AlertRule}, che fornisce il comportamento condiviso tra le diverse regole, in particolare per 
quanto riguarda l'accesso ai parametri configurati dall'utente e la gestione dei permessi di utilizzo.
\newline
Le singole regole concrete estendono quindi \texttt{AlertRule} e implementano la logica specifica necessaria per individuare le relative violazioni. Ogni regola definisce 
inoltre i metadati utilizzati dall'interfaccia utente e implementa principalmente due metodi:
\begin{itemize}
\item \texttt{check}: esegue le verifiche necessarie, effettuando le interrogazioni al database per individuare le condizioni che soddisfano i criteri della regola;
\item \texttt{toAlert}: elabora le informazioni relative alla violazione individuata e le converte in un \acrfull{dto}, che verrà successivamente utilizzato per la generazione 
  dell'alert.
\end{itemize}
Le regole implementate rispondono a diverse esigenze operative degli utenti finali di \textit{ErpAI}, con l'obiettivo di individuare tempestivamente situazioni 
che possono richiedere un intervento. Tra gli esempi implementati sono presenti \textit{LateBuyOrdersRule} e \textit{LateInvoicesRule}, che individuano rispettivamente 
ordini di acquisto e fatture in ritardo rispetto alle condizioni previste, \textit{SchedulingErrorsRule} e \textit{ForecastErrorsRule}, che segnalano errori o discrepanze nei processi di scheduling della 
produzione e nelle predizioni delle vendite, e \textit{PausedMachineRule}, che individua le macchine che non stanno eseguendo alcuna lavorazione, fornendo anche informazioni relative alla durata e alla causa 
della pausa.

\subsection{Jobs}
I \textit{job} sono le componenti responsabili dell'esecuzione asincrona delle operazioni necessarie alla generazione e alla distribuzione degli alert. Essi vengono avviati 
dal sistema di pianificazione ed eseguono i controlli definiti dalle regole, la persistenza degli alert e la successiva gestione delle notifiche.

\subsubsection*{AlertJob}
La classe \texttt{AlertJob} costituisce il principale punto di ingresso del motore di generazione degli alert. Essa espone i metodi \texttt{handle}, che contiene 
la logica principale di elaborazione, e \texttt{tooManyAlerts}, utilizzato per gestire il caso in cui il numero di violazioni individuate superi la soglia massima configurata.
\newline
Il metodo \texttt{handle} esegue le \textit{rules} abilitate dall'utente e raccoglie i risultati delle verifiche. Gli alert individuati vengono quindi persistiti nel database. 
Successivamente viene verificato il numero di alert che devono essere notificati all'utente. Se tale numero è inferiore alla soglia massima prevista, vengono create le relative 
notifiche e viene effettuato il \textit{dispatch} di un \texttt{BroadcastAlertJob} per la loro distribuzione all'interfaccia utente.
\newline
L'esecuzione della distribuzione delle notifiche tramite una coda asincrona permette di separare la generazione degli alert dalla loro visualizzazione, evitando che un numero 
elevato di notifiche rallenti l'elaborazione principale. La gestione delle code è affidata a Laravel Horizon, che utilizza Redis come sistema di coda.
\newline
Se invece il numero di alert supera la soglia configurata, viene richiamato il metodo \texttt{tooManyAlerts}. In questo caso il sistema evita la generazione di un numero 
eccessivo di notifiche e produce una notifica specifica che informa l'utente della presenza di un numero elevato di violazioni.

\subsubsection*{BroadcastAlertJob}
La classe \texttt{BroadcastAlertJob} si occupa della distribuzione degli alert verso l'interfaccia utente. Il job crea la \textit{Notification Filament} corrispondente all'alert 
e ne gestisce il \textit{broadcasting}, consentendo alla notifica di essere recapitata all'utente destinatario.
\newline
L'esecuzione di questa operazione tramite un job separato permette di mantenere indipendente la generazione degli alert dalla loro distribuzione, sfruttando il sistema di code 
asincrone fornito da Laravel Horizon e Redis.

\subsection{Command}
Il motore di generazione degli alert viene avviato tramite un comando Artisan, eseguibile manualmente da terminale mediante \texttt{php artisan alert} oppure 
automaticamente attraverso il sistema di scheduling di Laravel. In ambiente di produzione, il comando viene registrato nello scheduler dell'applicazione, definito in 
\texttt{bootstrap/app.php}, che ne pianifica l'esecuzione secondo intervalli temporali prestabiliti.

\subsubsection*{AlertCommand}
La classe \texttt{AlertCommand} svolge esclusivamente il ruolo di punto di ingresso del motore di generazione degli alert e non contiene direttamente la logica 
necessaria all'elaborazione delle regole.
\newline
Il metodo \texttt{handle}, unica operazione implementata dal comando, individua gli utenti che dispongono di almeno una regola alert abilitata e procede al \textit{dispatch} di 
un \texttt{AlertJob} per ciascuno di essi. L'elaborazione vera e propria viene quindi delegata ai job, permettendo al comando di terminare rapidamente e lasciando l'esecuzione 
dei controlli al sistema di code asincrone.

\subsection{Resources e Pages}
Per la realizzazione dell'interfaccia relativa al sistema di alert sono state sviluppate, utilizzando le funzionalità messe a disposizione da Filament, una 
\textit{Resource} e due \textit{Pages}. Queste componenti sono implementate all'interno dell'applicazione Laravel e permettono di integrare direttamente la logica del 
sistema di alert con gli elementi dell'interfaccia utente.
\begin{itemize}
  \item \textbf{\texttt{AlertResource}:} risorsa che gestisce la visualizzazione della lista degli alert rilevati;
  \item \textbf{\texttt{ViewAlert}:} pagina dedicata alla visualizzazione del dettaglio di un singolo alert. Fornisce inoltre le funzionalità necessarie per richiedere una 
    spiegazione dell'evento tramite il modulo XAI e per raggiungere la risorsa del gestionale associata all'evento che ha generato l'alert;
  \item \textbf{\texttt{AlertSettings}:} pagina dedicata alla configurazione del sistema di alert, attraverso la quale l'utente può aggiungere, configurare e rimuovere le regole 
    alert disponibili.
\end{itemize}

\section{Sistema di logging delle operazioni}
Il sistema di logging delle operazioni si basa sull'utilizzo di una classe \textit{Observer}, incaricata di intercettare le principali operazioni effettuate sui modelli 
Eloquent presenti nel sistema. In particolare, vengono registrati a database gli eventi relativi alla creazione, alla modifica e all'eliminazione dei record, permettendo 
di mantenere uno storico delle operazioni effettuate dagli utenti.

\subsection{Observers}

\subsubsection*{AuditObserver}
\texttt{AuditObserver} è la classe \textit{Observer} responsabile della registrazione delle operazioni effettuate sui modelli Eloquent. Essa implementa i 
metodi \texttt{created()}, \texttt{updated()} e \texttt{deleted()}, che vengono richiamati automaticamente dal framework in corrispondenza delle relative operazioni.
I metodi dell'\textit{Observer} delegano quindi la costruzione del record di audit a un metodo comune, \texttt{base}, che predispone un'istanza di \texttt{AuditLog} contenente 
le informazioni necessarie alla registrazione dell'operazione. Il record viene successivamente completato con i dati specifici dell'evento e persistito nel database.

\subsection{Resources e Pages}
Per la realizzazione dell'interfaccia relativa al sistema di logging sono state sviluppate una \textit{Resource} e una \textit{Page}, quest'ultima integrata con una vista 
\textit{Blade} dedicata alla visualizzazione dei dettagli.
\begin{itemize}
  \item \textbf{\texttt{AuditLogResource}:} risorsa che gestisce la visualizzazione dell'elenco delle operazioni registrate nel sistema;
  \item \textbf{\texttt{ViewAuditLogPage}:} pagina dedicata alla visualizzazione del dettaglio di un singolo log. A differenza della pagina \texttt{ViewAlert}, non utilizza 
    direttamente i componenti di visualizzazione dei record forniti da Filament, ma passa i dati alla relativa vista \textit{Blade} per la loro elaborazione;
  \item \textbf{\texttt{view-audit-log.blade}:} vista \textit{Blade} responsabile della visualizzazione delle modifiche effettuate sul modello associato al log. L'utente può 
    scegliere se visualizzare l'intero storico dell'istanza oppure limitarsi all'operazione selezionata. È inoltre possibile visualizzare tutte le proprietà del modello oppure 
    esclusivamente quelle interessate dalla modifica. Per ogni operazione vengono infine riportati la data e l'utente responsabile dell'azione.
\end{itemize}

\section{Design Pattern utilizzati}
\subsection{Observer Pattern}
La classe \texttt{AuditObserver} implementa l'\textit{Observer Pattern} e viene registrata all'interno di \texttt{AppServiceProvider}, classe responsabile della registrazione 
e configurazione dei servizi dell'applicazione. Per ogni modello Eloquent soggetto a tracciamento viene effettuata la registrazione dell'Observer tramite il metodo 
\texttt{\$model::observe(AuditObserver::class)}.
\newline
A seguito della registrazione, Laravel provvede a notificare automaticamente \texttt{AuditObserver} quando sul modello si verificano gli eventi \texttt{created}, 
\texttt{updated} o \texttt{deleted}. L'Observer intercetta quindi tali eventi e provvede alla creazione del relativo record di \texttt{AuditLog}, permettendo di mantenere 
una registrazione centralizzata delle operazioni effettuate sui modelli del sistema.

\subsection{Strategy Pattern}
Il \textit{Strategy Pattern} viene utilizzato per rendere intercambiabili le implementazioni dei provider LLM e delle regole alert. Entrambe le componenti sono definite 
attraverso un'interfaccia comune, che permette al resto dell'applicazione di interagire con le diverse implementazioni senza dipendere direttamente dalle loro caratteristiche 
specifiche.
\newline
Nel modulo XAI, \texttt{AiProviderContract} definisce le operazioni comuni che devono essere fornite dai diversi provider. Le classi concrete, come \texttt{RunpodOllamaProvider}, 
implementano tale contratto fornendo le modalità specifiche di comunicazione con il servizio LLM. In questo modo il provider utilizzato può essere sostituito senza modificare la 
logica del servizio che gestisce la conversazione.
\newline
Un approccio analogo viene adottato per il sistema di alert tramite \texttt{AlertRuleContract}. Ogni regola concreta implementa lo stesso contratto, fornendo una propria 
strategia per individuare le violazioni e trasformarle in alert. Il motore di generazione può quindi eseguire le diverse regole senza dover conoscere nel dettaglio la logica 
specifica di ciascuna di esse.
\newline
L'utilizzo di questo pattern permette di aggiungere nuovi provider o nuove regole alert mantenendo separata la loro logica specifica dal resto dell'applicazione e riducendo le 
modifiche necessarie alle componenti già esistenti.

\subsection{Command Pattern}
Il sistema di generazione degli alert utilizza un comando dedicato, \texttt{AlertCommand}, che costituisce il punto di ingresso per l'avvio dell'elaborazione. Il comando 
incapsula l'operazione di avvio del controllo delle regole alert e ne separa l'esecuzione dalla logica che implementa effettivamente il processo.
\newline
Il metodo \texttt{handle()} di \texttt{AlertCommand} individua gli utenti che possiedono regole alert abilitate e provvede al \textit{dispatch} di un \texttt{AlertJob} per 
ciascuno di essi. In questo modo il comando non contiene direttamente la logica di verifica delle regole, ma si limita ad avviare il processo, demandandone l'esecuzione ai job.
\newline
Il comando può essere eseguito manualmente tramite Artisan oppure automaticamente attraverso lo scheduler dell'applicazione. Questa separazione permette di mantenere indipendenti 
il meccanismo di avvio e la logica di elaborazione degli alert.

\subsection{Template Method Pattern}
Il \textit{Template Method Pattern} viene utilizzato nella gestione dei provider del modulo XAI. La classe astratta \texttt{AbstractAiProvider} definisce il 
comportamento comune necessario per elaborare una richiesta al modello linguistico, centralizzando le operazioni condivise tra i diversi provider.
\newline
In particolare, il metodo \texttt{reply()} gestisce il flusso generale dell'elaborazione, occupandosi della costruzione della conversazione, del caricamento dei tool, 
della gestione del contesto e dei parametri relativi alla richiesta. Le implementazioni concrete dei provider estendono quindi la classe astratta e forniscono le informazioni e 
i comportamenti specifici necessari per comunicare con il relativo servizio.
\newline
Lo stesso approccio viene utilizzato anche nella struttura dei tool e delle regole alert, attraverso le rispettive classi astratte \texttt{AbstractTool} e \texttt{AlertRule}, 
che centralizzano la gestione delle funzionalità comuni lasciando alle classi concrete l'implementazione della logica specifica.

\newpage